\documentclass[12pt, a4paper]{report}

% Pacotes essenciais para textos em português
\usepackage[utf8]{inputenc} % Codificação do texto
\usepackage[T1]{fontenc}    % Codificação de fontes
\usepackage[portuguese]{babel} % Idioma
\usepackage{graphicx}       % Para inserir imagens
\usepackage{hyperref}       % Para links clicáveis no sumário

% Add-ons 
% \usepackage[square]{natbib} % para utilizar o \citeauthor{}

% -----------------------------------------------------------------------------------------------
\begin{document}
% -----------------------------------------------------------------------------------------------

\cite{10.1007/978-3-642-19513-6_12}


\cite{digitalcorpora}



\cite{GARFINKEL2009S2} 


\cite{alherbawi2016survey}





File carving is a forensics technique used to recover files based solely on file structure and content, without relying on matching file system meta-data, and is most often employed to recover files from the unallocated space in a drive.

% ==========================================================
% Tabela: Mapeamento de artigos
% ==========================================================
\begin{table}[h]
    \caption{Mapeamento de Artigos e Utilidade na Pesquisa}
    \centering
    \begin{tabular}{p{5cm}|c|p{7cm}}
        \hline
        \textbf{Nome do Artigo / Tópico} & 
        \textbf{Citação} & 
        \textbf{Como pretendo usar / Como pode ser útil} \\
        \hline
        \hline
        Título do Artigo sobre Carving & \cite{10.1007/978-3-642-19513-6_12} & Utilizado para descrever os métodos X e Y na revisão da literatura. \\
        \hline
        Portal Digital Corpora & \cite{digitalcorpora} & Fonte principal para download das imagens de disco usadas nos testes práticos. \\
        \hline
        Bringing science to digital forensics... & \cite{GARFINKEL2009S2} & Base teórica para justificar a escolha de datasets padronizados no experimento. \\
        \hline
        A survey on file carving... & \cite{alherbawi2016survey} & Fornece o estado da arte e comparações de ferramentas que embasam o capítulo 2. \\
        \hline
    \end{tabular}
    % \usepackage{abntex2}        % Para o comando \legend
    % \legend{Fonte: O Autor (2026).}
    \label{tbl:mapeamento_referencias}
\end{table}

\newpage
\begin{table}[h]
\caption{Mapeamento de Datasets}
\centering
    \begin{tabular}{p{5cm}|c|p{7cm}}

    \hline
    \textbf{Nome do dataset} & 
    \textbf{Descrição} & 
    \textbf{link} \\  \hline
    \hline
    Dataset de imagens & 
    % \cite{pan2025pairhumanhighfidelityphotographicdataset} &
    Dataset de imagens hq & 
    \cite{pan2025pairhumanhighfidelityphotographicdataset}
                    

    

\end{tabular}
\label{tbl:mapeamento_datasets}
\end{table}

% ---------------------------------------------
% [Pacote]: Útil para usar o \citeauthor
% \usepackage[square]{natbib}
% \citeauthor{}

% ================================================
\bibliographystyle{abntex2-alf}
\bibliography{biblio}

\end{document}